% ============================================================
% Amazon-VT CFP 2026-2027 — 1-Page Abstract v3
% Provenance-Aware Agentic Knowledge Graph Systems
% Format: proposalnsf
% ============================================================

\documentclass{files/proposalnsf}

\usepackage[utf8]{inputenc}
\usepackage[T1]{fontenc}
\usepackage{graphicx}
\usepackage{hyperref}
\usepackage{microtype}
\usepackage{cite}
\usepackage{enumitem}
\setlist{nosep, leftmargin=1.5em, topsep=0.2em}

\usepackage{xspace}
\usepackage{xcolor}

\newcommand{\todo}[1]{{\color{red}\bfseries [[TODO: #1]]}}

\newcommand{\tool}{\textcolor{blue}{\texttt{cool\_tool}}\xspace}

\newcommand{\approach}{\textcolor{blue}{\textit{approach}}\xspace}

\newcommand{\proposalEleven}{\textbf{Provenance-Aware Agentic Knowledge Graphs}}
\newcommand{\proposalEighteen}{\textbf{Provenance-Aware Agentic Knowledge Graphs}}
\newcommand{\proposalNineteen}{\textbf{Provenance-Aware Agentic Knowledge Graphs}}


\newcommand{\ie}{\textit{i.e.,}\xspace}

\newcommand{\eg}{\textit{e.g.,}\xspace}

\newcommand{\etal}{\textit{et al.}\xspace}

\begin{document}

\begin{center}
    \large{\proposalNineteen} \\ \vspace{0.5pt}
\footnotesize{PI: Chris Brown, Computer Science, Virginia Tech} \\ 
\footnotesize{Co-PI: Xuan Wang, Computer Science, Virginia Tech} \\
\footnotesize{Co-PI: Scott T. Steinmetz, Sandia National Labs}
\end{center}

\small{\noindent\textbf{Focus Area:} 19 -- AI for Scientific Reasoning} \hfill \\
\small{\textbf{Topic:} B. Hypothesis Generation from Multi-Modal Data}

\subsubsection*{Project Summary / Abstract}

Generating novel, testable scientific hypotheses from multi-modal data requires synthesizing knowledge across heterogeneous sources---scientific literature, experimental datasets, simulation outputs---while maintaining traceable links between each hypothesis and the specific evidence supporting it. Large language models can fluently generate hypotheses, but cannot enforce \textit{span-level provenance}: the traceable chain from a generated claim back to the source passages, entities, and relations that license it~\cite{Ji2023HallucinationSurvey}. Without such provenance, hypothesis generation produces plausible-sounding text rather than verifiable scientific inference---a fundamental barrier to AI as a true partner in the scientific method~\cite{Asai2024SelfRAG}.

Our prior work identifies three interdependent gaps that jointly prevent AI-assisted hypothesis generation from achieving scientific rigor (Cusati \& Brown, ICSE-FoSE 2026~\cite{cusati2026fose}): (1)~an \textit{architectural capability gap}---no existing system integrates persistent structured knowledge graphs (KGs) with autonomous multi-agent research workflows; (2)~an \textit{evaluation gap}---no validated framework measures the quality of AI-generated research outputs at the level of evidence grounding and provenance completeness; and (3)~a \textit{provenance gap}---no generalizable mechanism enforces span-level evidence alignment as an architectural invariant across agent operations. These gaps form a cycle: without persistent knowledge infrastructure, provenance is infeasible; without provenance, evaluation is unreliable; without evaluation, architecture design lacks a target to optimize against.

This project proposes \textbf{PA-AKG (Provenance-Aware Agentic Knowledge Graphs)}, a provenance-first architecture that breaks this cycle by making provenance enforcement \textit{architectural}, not post-hoc. PA-AKG organizes around three layers: (1)~a \textbf{knowledge representation layer}---a persistent research KG where problems, datasets, methods, constraints, and evidence are first-class entities linked by explicit semantic relations~\cite{Hogan2021KnowledgeGraphs}; (2)~an \textbf{automation layer}---span-level evidence alignment (automated claim extraction → source passage linkage) and canonicalization auditing (entity resolution to canonical KG nodes), building structured claim--evidence pairs as the substrate for hypothesis generation; and (3)~an \textbf{agentic orchestration layer}---ranking agents (prioritize research directions by tractability and evidence strength), continuation agents (propose follow-on analyses based on KG gaps), evaluation agents (execute reproducible workflows), and synthesis agents (write back structured artifacts). The key invariant: every agent handoff carries provenance constraints requiring downstream agents to inherit and extend evidence chains.

This architecture directly implements DOE 19B's vision: ``AI systems capable of synthesizing information from diverse sources to identify gaps in current knowledge and formulate novel, testable hypotheses.'' Crucially, PA-AKG produces \textit{traceable} hypotheses---each is accompanied by the specific evidence chain that supports it, the KG nodes it derives from, and the span-level citations linking it back to source documents---making AI-generated hypotheses auditable, reproducible, and refutable.

Expected outcomes include: the PA-AKG architecture and open prototype; evaluation framework E1--E4 (span-level alignment precision/recall, canonicalization consistency, traceability completeness, longitudinal stability); empirical benchmarks comparing PA-AKG against vector RAG, GraphRAG~\cite{Edge2024GraphRAG}, and static snapshot baselines; and a generalizable provenance-first infrastructure for traceable AI-assisted scientific discovery.

\subsubsection*{Research Plan}

\paragraph{Introduction.}
DOE's 19B challenge calls for AI systems capable of synthesizing multi-modal scientific knowledge to identify gaps and formulate testable hypotheses. Existing approaches fail at the fundamental requirement: they generate outputs without enforcing that each claim is traceable to specific supporting evidence. Retrieval-augmented generation~\cite{Lewis2020RAG} improves factual grounding but does not maintain persistent, structured knowledge across sessions, cannot reason over relations between entities, and provides no mechanism for auditing hypothesis derivations. Graph-augmented approaches~\cite{Edge2024GraphRAG} add graph structure but do not enforce provenance as an architectural constraint---provenance is a reporting artifact rather than a prerequisite for agent operations.

Our prior work establishes both the need and the feasibility of PA-AKG. Brown \& Cusati (ESEM 2025~\cite{brown2025esem}) empirically demonstrate that generative AI tools propagate evidence inconsistently across SE contexts, producing outputs that cite sources without verifying that those sources support the specific claims made. Cusati \& Brown (ICSE-FoSE 2026~\cite{cusati2026fose}) demonstrate that research knowledge in software engineering accumulates in structurally disconnected artifacts---individual papers, data files, tool outputs---without persistent structured infrastructure enabling AI to reason over the \textit{relationships} between them. Together, these studies establish that the hypothesis generation problem is fundamentally a knowledge infrastructure problem, and that solving it requires PA-AKG's provenance-first architecture.

\paragraph{Objectives.}
\begin{enumerate}
    \item Design and implement the PA-AKG three-layer architecture with provenance enforcement as an architectural invariant across all agent operations.
    \item Implement span-level evidence alignment: automated claim extraction producing structured claim--evidence pairs with source passage linkage.
    \item Implement canonicalization auditing: entity normalization, validation, and deduplication producing canonical KG nodes with inter-rater-reliable entity resolution.
    \item Build multi-agent orchestration with ranking, continuation, evaluation, and synthesis agents operating under provenance inheritance constraints.
    \item Evaluate PA-AKG using the E1--E4 evaluation framework against vector RAG, GraphRAG, and static snapshot baselines on SE research hypothesis generation tasks.
\end{enumerate}

\paragraph{Research Questions.}
\begin{itemize}
    \item \textbf{RQ1:} How accurately does automated span-level evidence alignment link generated claims to source passages compared to human-annotated gold spans (precision/recall, E1)?
    \item \textbf{RQ2:} How consistently does canonicalization auditing resolve entities to canonical KG nodes relative to expert entity clusters (Cohen's $\kappa$, E2)?
    \item \textbf{RQ3:} What fraction of KG nodes maintain complete provenance chains as the research corpus evolves over time (traceability completeness, E3)?
    \item \textbf{RQ4:} How does PA-AKG's structured provenance enforcement affect hypothesis quality, specificity, and falsifiability compared to vector RAG and GraphRAG baselines?
\end{itemize}

\paragraph{Detailed Research Plan.}

\textit{Phase 1 (Months 1--3): KG Architecture and Evidence Alignment Pipeline.}
We will design the research KG schema with first-class entity types: \textsc{Problem}, \textsc{Dataset}, \textsc{Method}, \textsc{Constraint}, \textsc{Evidence}, and \textsc{Hypothesis}, linked by semantic relations (\textsc{supports}, \textsc{contradicts}, \textsc{enables}, \textsc{requires}). We will implement the claim extraction pipeline using structured prompting to produce claim--evidence pairs, followed by span-level alignment that links each extracted claim to its supporting source passages with character-level span annotations. We will implement canonicalization auditing: entity normalization via embedding-based clustering with deduplication, followed by validation against KG-encoded entity constraints. We will build the hybrid retrieval layer combining a property graph store (Neo4j or equivalent) with a vector index, enabling both structured relational queries and approximate semantic similarity search. Primary validation corpus: the software engineering research literature corpus from Cusati \& Brown prior work~\cite{cusati2026fose}, which provides a bounded, well-understood domain for initial evaluation.

\textit{Phase 2 (Months 4--6): Agentic Orchestration with Provenance Inheritance.}
We will implement the four PA-AKG agent types. \textit{Ranking agents} score candidate research directions using a composite tractability score combining evidence strength (number of supporting spans), coverage (fraction of KG entities referenced), and novelty (distance from existing KG hypotheses). \textit{Continuation agents} propose follow-on analyses by querying the KG for under-evidenced entities and suggesting experiments or literature searches to fill identified gaps. \textit{Evaluation agents} execute reproducible experimental workflows by invoking domain-specific analysis tools and writing structured results back to the KG as evidence nodes with VVUQ-style verification receipts. \textit{Synthesis agents} generate structured hypothesis artifacts (problem statement, supporting evidence chain, proposed experimental test, predicted outcome) and write them back to the KG as \textsc{Hypothesis} nodes with full provenance metadata. The provenance inheritance constraint is enforced at agent handoff: each agent receives a provenance context object carrying accumulated evidence chain references, and may only produce outputs that extend---not discard---this chain.

\textit{Phase 3 (Months 7--9): Evaluation and Generalization.}
We will evaluate PA-AKG using the E1--E4 framework. \textbf{E1} (span-level alignment): compute precision/recall of PA-AKG span annotations against a human-annotated gold set of 500 claim--evidence pairs from the SE research corpus. \textbf{E2} (canonicalization consistency): measure Cohen's $\kappa$ between PA-AKG entity cluster assignments and expert-produced canonical entity clusters. \textbf{E3} (traceability completeness): measure percentage of KG nodes with complete provenance chains (all required evidence links present) as a function of corpus size and evolution rate. \textbf{E4} (persistence): track longitudinal stability score (fraction of hypotheses remaining valid as corpus evolves) across six simulated corpus update cycles. We will compare against four baselines: (B1) vector RAG without span alignment; (B2) GraphRAG~\cite{Edge2024GraphRAG} without provenance enforcement; (B3) static snapshot system (no cumulative KG); and (B4) LLM-only generation without retrieval. Secondary domain validation will extend the PA-AKG pipeline to one additional research domain to test generalizability of the schema and alignment pipeline.

\paragraph{Team Roles.}
\textbf{PI Chris Brown} leads project direction, KG architecture design, evaluation methodology (E1--E4), and SE-domain corpus curation. \textbf{Co-PI Xuan Wang} leads multi-agent orchestration design, ranking and continuation agent implementation, and embedding-based canonicalization. \textbf{Co-PI Scott T. Steinmetz} (Sandia National Labs) leads provenance infrastructure design, formal verification of evidence chains, and secondary domain validation against DOE science corpora. \textbf{GRA Jason Cusati} implements the PA-AKG system end-to-end, constructs and annotates the evaluation corpus, and executes experimental evaluation and baseline comparisons.

\bibliographystyle{IEEEtran}
\bibliography{references}

\end{document}
